\chapter{El Modelo Operativo y las Proyecciones de una Startup}\label{ch:operating-model}

% Alcance: puente de economía unitaria a P&L (construcción de ingresos por cohorte), el modelo
%   de tres palancas para EBITDA, presupuesto por escenarios y análisis de sensibilidad, runway.
%   Capítulo nuevo (cierra la brecha en modelos financieros).
% Conecta la economía unitaria (\cref{ch:customer-economics}) con el DCF (\cref{ch:valuation}):
%   construye las proyecciones de FCF que el DCF da por sentadas.
% Referencias cruzadas: \cref{ch:customer-economics} (CLV/CAC), \cref{ch:valuation} (DCF),
%   \cref{ch:fundraising-and-dilution} (disparador de runway), \cref{ch:financial-statements} (tres estados).
% Borrador según estilo de la casa: prosa continua + separadores example/admonition;
%   ejemplos trabajados con la SaaS de referencia.

\section{El Puente de Economía Unitaria al P\&L}\label{sec:unit-econ-to-pl}
% complexity: 3

¿Cómo se convierte un conjunto de supuestos de economía unitaria en un pronóstico de ingresos?
La respuesta es la pila de cohortes: los nuevos clientes de cada mes entran en una cohorte que
se reduce por deserción en cada mes posterior, contribuyendo al ARR mientras permanezcan
activos. Al sumar las contribuciones supervivientes de cada cohorte se obtiene el ARR actual de
la empresa. No se trata de un ejercicio descendente de participación en el TAM; es una identidad
contable ascendente que hace que toda proyección de ingresos sea trazable a dos palancas
operativas --- entrada de nuevos clientes y \emph{retention} --- y una palanca de precio, el ARPU.

Sea $N_s$ el número de nuevos clientes adquiridos en el mes $s$, ARPU el ingreso anualizado por
cliente y $c$ la tasa mensual de \emph{churn}. El ARR al final del período $t$ es la suma sobre
todas las cohortes anteriores de los clientes aún activos:
\begin{equation}\label{eq:cohort-revenue}
  \mathrm{ARR}_{t}
  \;=\;
  \mathrm{ARPU}
  \sum_{s=1}^{t}
  N_{s}\,(1-c)^{\,t-s}.
\end{equation}
Cuando $N_s$ es constante en $N$ y la serie lleva el tiempo suficiente para aproximarse al estado
estacionario, la suma colapsa a $N\,\mathrm{ARPU}/c$; el techo de estado estacionario exacto se
alcanza sólo asintóticamente, por lo que las empresas reales se encuentran siempre en algún
punto de la curva de crecimiento por debajo de él.
\Cref{eq:cohort-revenue} es el gemelo de ingresos de la fórmula de CLV: \cref{eq:clv} valora a
un único cliente, mientras que \cref{eq:cohort-revenue} apila la misma mecánica de supervivencia
en todas las cohortes del libro. La identidad de la tasa de crecimiento de MRR en
\cref{eq:mrr-growth} es la derivada de \cref{eq:cohort-revenue} con respecto al tiempo.

\begin{example}[Pila de cohortes hacia ARR y FCF del Año~1]
La empresa adquiere \num{200} nuevos clientes por mes a ARPU \money{600}/año, con \emph{churn}
mensual de \qty{0.65}{\percent}. Tras \num{38} meses de rampa, \cref{eq:cohort-revenue} produce:
\[
  \mathrm{ARR}_{38}
  \;=\; \money{600}\times
        \frac{200\,\bigl(1-0.9935^{38}\bigr)}{0.0065}
  \;\approx\; \money{4050000},
\]
lo que redondea al ARR de \money{4000000} del plan operativo. Con esta base, capital invertido
de \money{3000000} y ROIC del \qty{20}{\percent} (\cref{eq:roic}), el NOPAT equivale a
\(\money{3000000}\times 0.20=\money{600000}\). Dado que una empresa SaaS tiene gasto de capital
casi nulo y variaciones mínimas en el capital de trabajo, el flujo de caja libre es
aproximadamente igual al NOPAT; \emph{el FCF del Año~1 de \money{600000}} alimenta directamente
el DCF en \cref{eq:dcf} como ancla del período de pronóstico explícito, cuya suma de valores
presentes es aproximadamente \money{3377825}.

La implicación es operativa: tanto el ARR de \money{4000000} como el FCF de \money{600000} son
resultados derivados de tres tasas observables --- entrada, \emph{retention} y precio. Un
fundador que no pueda nombrar el valor actual de las tres no puede cerrar el ciclo desde el
gasto en \emph{marketing} hasta el valor de la empresa.
\end{example}

\begin{admonition}{Advertencia}
Las estadísticas de \emph{retention} combinadas ocultan la heterogeneidad entre cohortes. Una
empresa que adquiere la mitad de sus clientes a través de un canal de baja \emph{retention}
reportará un \emph{churn} agregado aceptable mientras las cohortes saludables subvencionan a las
deficientes. Modele la \emph{retention} \emph{por cohorte de adquisición}; la mezcla en la cima
del \emph{funnel} es tan importante como la tasa en sí.
\end{admonition}

La pila de cohortes se conecta directamente al P\&L: el ARR es la línea de ingresos; la
estructura de costos variables de \cref{ch:financial-statements} --- margen bruto del
\qty{70}{\percent}, margen de contribución del \qty{42}{\percent} --- determina qué fracción
del ingreso de cada cohorte fluye hacia el EBITDA. La siguiente sección convierte esas tasas en
un modelo de palancas explícito.


\section{El Modelo de Tres Palancas}\label{sec:three-driver-model}
% complexity: 3

Una vez que el ARR se proyecta a partir de la pila de cohortes, el P\&L se deriva de tres
tasas: el margen bruto, la base de costos fijos y su composición entre los tres rubros de gastos
operativos. El margen bruto refleja la ingeniería del producto y no puede modificarse
rápidamente; la composición de costos refleja decisiones de gestión sobre dónde invertir. La
línea de EBITDA no es más que el residuo aritmético de aplicar estas tasas a una base de ARR en
crecimiento. La disciplina de expresar el modelo de esta forma obliga al equipo a ser explícito
sobre el momento en que el apalancamiento operativo se activa --- el punto en que el crecimiento
del ARR supera al de la base de costos fijos y el EBITDA se vuelve positivo.

Formalice la relación como:
\begin{equation}\label{eq:op-model}
  \mathrm{EBITDA}_{t}
  \;=\;
  \mathrm{ARR}_{t}\cdot g_{\mathrm{m}}
  \;-\;
  \mathrm{FixedCosts}_{t},
\end{equation}
donde $g_{\mathrm{m}}$ es la tasa de margen bruto y $\mathrm{FixedCosts}_{t}$ es la suma de
Ventas y \emph{Marketing} (S\&M), Investigación y Desarrollo (R\&D) y Gastos Generales y
Administrativos (G\&A). En una empresa SaaS el margen bruto es relativamente estable (el costo
de atender a un cliente adicional a escala es cercano a cero), de modo que la expansión del
EBITDA depende principalmente de la rapidez con que el crecimiento del ARR supere al de los
costos fijos --- la definición canónica del apalancamiento operativo.
\textcite{koller2025valuation} rastrean esta dinámica a través del puente EBITDA; \textcite{ross2022}
ofrecen el tratamiento de finanzas corporativas del apalancamiento operativo y sus implicaciones
para el riesgo del capital.

\begin{example}[Progresión del modelo de tres palancas: margen EBITDA y Regla del 40]
Partiendo de un ARR de \money{4000000} que crece al \qty{30}{\percent} anual, con margen bruto
del \qty{70}{\percent}, la mezcla de gastos operativos evoluciona conforme la empresa escala:

\medskip
\noindent\resizebox{\linewidth}{!}{%
\begin{tabular}{lrrrr}
\toprule
 & Año~0 & Año~1 & Año~2 & Año~3 \\
\midrule
ARR                       & \money{3077000} & \money{4000000} & \money{5200000} & \money{6760000} \\
Beneficio bruto (70\%)    & \money{2154000} & \money{2800000} & \money{3640000} & \money{4732000} \\
S\&M (\% del ARR)         & 50\%       & 50\%       & 45\%       & 35\%       \\
R\&D (\% del ARR)         & 25\%       & 25\%       & 25\%       & 20\%       \\
G\&A (\% del ARR)         & 10\%       & 10\%       & 15\%\textsuperscript{*} & 8\%        \\
Total opex                & 85\%       & 85\%       & 85\%       & 63\%       \\
Margen EBITDA             & \(-15\%\)  & \(-15\%\)  & \;\;0\%    & \;\,+7\%   \\
\bottomrule
\end{tabular}%
}

\medskip
\noindent\textsuperscript{*}G\&A sube temporalmente en el Año~2 mientras la empresa invierte
en infraestructura legal y financiera antes de una ronda de financiamiento; se comprime al
\qty{8}{\percent} en el Año~3 a medida que el ARR escala por encima de la función escalón de
G\&A.

\medskip
La puntuación de la Regla del 40 (\cref{eq:rule40}) es la suma de la tasa de crecimiento del
ARR y el margen EBITDA: el Año~1 obtiene \(30+(-15)=15\), el Año~2 obtiene \(30+0=30\) y el
Año~3 obtiene \(30+7=37\) --- una progresión desde la baja rentabilidad profunda hasta acercarse
al umbral de referencia. La implicación: el relato del apalancamiento operativo sólo es creíble
si las curvas de G\&A y R\&D se aplanan genuinamente conforme el ARR escala. Cualquier modelo
que mantenga los porcentajes de opex constantes mientras el ARR crece no está modelando
apalancamiento operativo; está modelando un negocio de costo más margen.
\end{example}

\begin{admonition}{Advertencia}
El apalancamiento operativo actúa en ambas direcciones. La misma base de costos fijos que
produce expansión del EBITDA con un crecimiento del \qty{30}{\percent} genera pérdidas
aceleradas si el crecimiento cae al \qty{15}{\percent}. Someta el modelo a una prueba de
estrés al \qty{50}{\percent} del crecimiento planificado antes de presentarlo a un consejo
directivo: con un crecimiento de ARR del \qty{15}{\percent}, el margen EBITDA del Año~3 en el
ejemplo anterior permanece cerca del \(-5\%\) y la puntuación de la Regla del 40 no cruza el
umbral. Un plan que requiere ejecución perfecta para evitar una crisis de liquidez es un plan
sin margen para el error.
\end{admonition}

\Cref{eq:op-model} produce el EBITDA que alimenta el puente EBITDA (\cref{eq:nopat-bridge})
y en última instancia el flujo de caja libre que \cref{eq:dcf} descuenta. La siguiente sección
pregunta qué tan sensible es ese resultado a los supuestos que lo sustentan.


\section{Presupuesto por Escenarios y Análisis de Sensibilidad}\label{sec:scenario-budget}
% complexity: 4

Un modelo operativo de un solo escenario es un pronóstico disfrazado. La versión honesta del
mismo trabajo es una superficie de sensibilidad que mapea el resultado del modelo a lo largo de
un rango de supuestos de entrada, para que la dirección sepa qué palanca importa más y en qué
medida. Las dos palancas más peligrosas en un modelo de economía unitaria SaaS son CAC y
\emph{churn}: el CAC governa el costo de llenar la cima de la pila de cohortes; el \emph{churn}
governa la velocidad a la que se vacía. Ambos pueden moverse de forma independiente y su
interacción no es lineal.

La derivada parcial del EBITDA con respecto al CAC (manteniendo todo lo demás constante) es
simplemente el número de nuevos clientes adquiridos por período. Con \num{200} nuevos clientes
por mes:
\begin{equation}\label{eq:sensitivity}
  \begin{aligned}
    \frac{\partial\,\mathrm{EBITDA}}{\partial\,\mathrm{CAC}} \;=\; -N
    \quad&\Longrightarrow\quad \Delta\mathrm{CAC}=\money{100} \\
    &\Rightarrow\; \Delta\mathrm{EBITDA} = -\money{100}\times 200 = -\money{20000}\text{ por mes},
  \end{aligned}
\end{equation}
o \money{240000} por año. El signo negativo es fundamental: el CAC es una salida de efectivo
que impacta el EBITDA antes de que se realice un solo dólar de valor de vida. El \emph{churn}
opera a través de la línea de ARR mediante \cref{eq:cohort-revenue}: un aumento de la tasa de
\emph{churn} de \qty{0.35}{\percent} mensual (de \qty{0.65}{\percent} a \qty{1.00}{\percent})
reduce el recuento de clientes en estado estacionario en aproximadamente un \qty{35}{\percent},
lo que, con un ARR de \money{4000000}, representa una reducción estructural de \money{1400000}
en el techo de ingresos --- un orden de magnitud mayor que un año de exceso de gasto en CAC.

\begin{figure}[htbp]
  \centering
  \includegraphics[width=0.92\textwidth]{op-model-bridge}
  \caption{Cascada del modelo operativo: desde los insumos de economía unitaria (ARPU, margen
    bruto, \emph{churn}, nuevos clientes por mes) hasta el ARR y el EBITDA, y de ahí al flujo
    de caja libre. Las alturas de las barras muestran cómo cada palanca contribuye o sustrae;
    la barra más a la derecha es el FCF del Año~1 de \money{600000}.}
  \label{fig:op-model-bridge}
\end{figure}

\begin{example}[Análisis de sensibilidad bidireccional: NPV a 3 años en función de CAC y churn]
La tabla siguiente mantiene ARR\textsubscript{0} en \money{4000000}, la adquisición anual de
nuevos clientes en \num{2400}, R\&D\,+\,G\&A en \money{2000000}, el margen bruto en
\qty{70}{\percent} y la tasa de descuento en \qty{11}{\percent}. Cada celda es el valor
presente neto a tres años del NOPAT, con una tasa impositiva del \qty{30}{\percent} aplicada
al EBITDA positivo.

\medskip
\noindent\resizebox{\linewidth}{!}{%
\begin{tabular}{lrrr}
\toprule
\multicolumn{1}{c}{CAC}
  & \multicolumn{1}{c}{Churn 0.50\%/mes}
  & \multicolumn{1}{c}{Churn 0.75\%/mes}
  & \multicolumn{1}{c}{Churn 1.00\%/mes} \\
\midrule
\money{400} & \money{2410000}  & \money{2120000}  & \money{1840000} \\
\money{500} & \money{2000000}  & \money{1710000}  & \money{1430000} \\
\money{600} & \money{1590000}  & \money{1300000}  & \money{1020000} \\
\money{700} & \money{1170000}  & \money{860000}   & \money{560000}  \\
\money{800} & \money{700000}   & \money{380000}   & \money{90000}   \\
\bottomrule
\end{tabular}%
}

\medskip
El caso base (CAC \money{600}, \emph{churn} \qty{0.65}{\percent}) produce un NPV a tres años
cercano a \money{1400000}, coherente con el VP del período explícito de \money{3377825} en el
DCF completo a cinco años (\cref{eq:scenario-ev}). El caso bajista --- CAC \money{800} y
\emph{churn} mensual del \qty{1.00}{\percent} --- reduce el NPV a apenas \money{90000}, un
colapso del \qty{94}{\percent} frente al caso base, ilustrando que las dos palancas se
componen: un CAC más alto eleva el lado de los costos justo en el momento en que un
\emph{churn} más alto comprime la base de ingresos. Nótese que ninguna celda de la tabla es
negativa; incluso el caso bajista genera FCF acumulado positivo en tres años. El riesgo en
\money{800}/\qty{1.00}{\percent} no es, por tanto, la insolvencia en este horizonte temporal,
sino un FCF insuficiente para justificar el múltiplo de valoración requerido para una ronda de
financiamiento.
\end{example}

\begin{admonition}{Advertencia}
Un modelo sensible a una sola palanca está mal calibrado. Si variar el CAC en un \qty{30}{\percent}
modifica el NPV en un \qty{50}{\percent} mientras que variar el \emph{churn} en la misma fracción
lo cambia sólo en un \qty{5}{\percent}, uno de los dos insumos no está siendo modelado con
suficiente granularidad. El análisis de sensibilidad no es sólo una herramienta de comunicación
de riesgos; es un diagnóstico que revela si la estructura del modelo coincide con las palancas
reales del negocio.
\end{admonition}

La superficie de sensibilidad alimenta la lógica de escenarios. Tres escenarios internamente
coherentes --- bajista, base y alcista --- deben contener cada uno su propia trayectoria de
EBITDA y su propio \emph{runway} implícito, porque la pregunta que el consejo directivo hace
siempre después de ver la tabla de sensibilidad es la misma: ¿cuánto tiempo falta para que
necesitemos recaudar capital?


\section{El \emph{Runway} y la Decisión de Financiamiento}\label{sec:runway}
% complexity: 2

El \emph{runway} es el puente entre el modelo operativo y la decisión de financiamiento.
Convierte la tasa neta de consumo de efectivo en un contador regresivo y establece la fecha
disparadora para la recaudación de capital. La mecánica es sencilla; la disciplina consiste en
ejecutar dos versiones de forma simultánea --- un \emph{runway} de crecimiento que asume que la
adquisición planificada de nuevos clientes continúa, y un \emph{runway} conservador que asume
crecimiento cero a partir de la base actual.

Sea $B$ el saldo de efectivo actual, $E$ el gasto operativo mensual (tasa de quema bruta o
\emph{burn} bruto) y $R_{\mathrm{m}}$ el efectivo cobrado mensualmente de los clientes (MRR
cobrado, no diferido):
\begin{equation}\label{eq:runway}
  \text{Runway (meses)}
  \;=\;
  \frac{B}{\,E - R_{\mathrm{m}}\,}
  \;=\;
  \frac{B}{\text{tasa de quema neta mensual}}.
\end{equation}
El denominador es el \emph{burn} neto: lo que pierde la cuenta bancaria cada mes después de
contabilizar las entradas de efectivo. La distinción entre \emph{burn} bruto y neto importa
porque una empresa de alto crecimiento con MRR sustancial tiene un reloj de supervivencia muy
diferente al que podría sugerir su estado de resultados. El \emph{burn} bruto es el piso
mínimo de costos; el \emph{burn} neto es la métrica existencial.
\textcite{brealey2020} desarrollan la mecánica del estado de flujos de efectivo que subyace a la
relación entre ingreso neto, cargos no monetarios y movimiento real de caja --- la misma
mecánica que explica por qué el MRR cobrado puede diferir del ingreso reconocido en el mes en
que se registra.

\begin{example}[Runway y disparador de recaudación de capital]
La empresa mantiene \money{2500000} en efectivo. El gasto operativo mensual (salarios,
infraestructura en la nube, salidas de S\&M) es de \money{650000}; el MRR actualmente cobrado
es de \money{340000}, coherente con la base de ARR de \money{4000000} y el hecho de que los
suscriptores anuales pagan mensualmente. Aplicando \cref{eq:runway}:
\[
  \text{Burn neto}
  \;=\; \money{650000} - \money{340000}
  \;=\; \money{310000}\text{ por mes},
\]
\[
  \text{Runway}
  \;=\; \frac{\money{2500000}}{\money{310000}}
  \;\approx\; \text{8 meses}.
\]
El \emph{disparador de recaudación} es la fecha que deja suficiente \emph{runway} para cerrar
una ronda. Un proceso competitivo de Serie A tarda de cuatro a seis meses desde el primer
contacto hasta la transferencia de fondos; agregar dos meses de margen establece el disparador
en seis meses de \emph{runway} restante. Ocho meses de \emph{runway} menos seis meses de plazo
necesario significa que la empresa debe comenzar la recaudación \emph{de inmediato} --- no en
la marca de los seis meses.

La implicación para el modelo operativo: el cálculo del \emph{runway} debe re-ejecutarse
mensualmente contra el efectivo real y el \emph{burn} neto real. Una tasa de \emph{burn} que
supera el plan en \money{50000} al mes durante tres meses ha consumido una ronda completa de
margen sin que aparezca ninguna línea en el estado de resultados. El capítulo sobre
financiamiento y dilución (\cref{ch:fundraising-and-dilution}) toma esta fecha disparadora y
trabaja la mecánica de valoración y las implicaciones del term sheet cuando se recauda capital
desde esta posición.
\end{example}

\begin{admonition}{Advertencia}
Ejecute siempre tanto el \emph{runway} de crecimiento como el conservador, y use el conservador
para establecer la fecha disparadora. El \emph{runway} de crecimiento es el escenario alcista;
el \emph{runway} conservador --- que mantiene el MRR plano al nivel actual e ignora las ventas
nuevas --- es el reloj de supervivencia. Una empresa que basa el calendario de su recaudación
en el \emph{runway} de crecimiento está apostando su supervivencia a que el plan se ejecute
exactamente como fue modelado. El proceso de recaudación no esperará una recuperación.
\end{admonition}

El modelo operativo está ahora completo como instrumento de planificación: \cref{eq:cohort-revenue}
construye el ARR a partir de insumos de economía unitaria observables; \cref{eq:op-model}
convierte el ARR y la estructura de costos en EBITDA; \cref{eq:sensitivity} cuantifica el
impacto de los cambios en palancas individuales; y \cref{eq:runway} convierte el perfil de
efectivo resultante en una línea temporal de supervivencia. El resultado de FCF de esta cadena
--- \money{600000} en el Año~1, creciendo a las tasas de \cref{eq:op-model} --- es el insumo
que el DCF en \cref{eq:dcf} descuenta hacia un valor empresarial.
\textcite{koller2025valuation} cierran el ciclo entre los resultados del modelo operativo y el
valor empresarial del DCF en los capítulos sobre pronósticos y palancas de valor.
